% ══════════════════════════════════════════════════════════════════════════════
\chapter{Préparation du corpus}
\label{chap:td1}
% ══════════════════════════════════════════════════════════════════════════════

La première étape du projet est la constitution du corpus documentaire au sens du
cours (chapitre~2, Vocabulaire). Un corpus est l'ensemble des textes sur lesquels portera
l'indexation ; sa qualité conditionne directement les performances de tout le système aval.
L'objectif est de transformer 326 fichiers HTML hétérogènes en un unique fichier
\texttt{corpus.xml} structuré, dont le schéma est imposé par le sujet. Ce fichier
servira de point d'entrée à toutes les étapes suivantes.

% ─────────────────────────────────────────────────────────────────────────────
\section{Objectif du TD}
% ─────────────────────────────────────────────────────────────────────────────

Ce TD constitue la première étape du projet LO17, dont le but est de construire un
système complet d'indexation et de recherche d'information sur une archive de bulletins
de veille technologique de l'ADIT. L'archive étudiée couvre plus de 300 articles extraits
de bulletins de veille sur la France entre 2011 et 2014. Chaque article contient des
méta-informations structurées (numéro, date, rubrique, auteur) ainsi que des contenus
(texte, images, contacts).

L'objectif concret est de préparer ce corpus en vue de son indexation future : à partir
des fichiers HTML bruts, produire un unique fichier \texttt{corpus.xml} structuré qui
rassemble l'ensemble des articles dans un format facilement exploitable. La méthode est
incrémentale — traiter un champ à la fois, valider sur plusieurs fichiers, puis généraliser
au corpus entier. La robustesse est une exigence explicite : toute absence de champ doit
être tolérée et journalisée sans interrompre le traitement.

% ─────────────────────────────────────────────────────────────────────────────
\section{Présentation du corpus ADIT}
% ─────────────────────────────────────────────────────────────────────────────

Les bulletins de l'ADIT (Agence pour la Diffusion de l'Information Technologique)
constituent une archive de veille technologique couvrant des thématiques variées :
biotechnologie, systèmes embarqués, matériaux, numérique. Chaque bulletin contient
plusieurs articles, chacun structuré en champs distincts (numéro, date, rubrique,
auteur, texte, images, contacts). Le corpus présente cependant plusieurs difficultés :
encodage déclaré en ISO-8859-1 alors que le contenu est en UTF-8, CSS classes non
documentées, navigation par tableaux HTML imbriqués, et présence de champs parfois absents.

Le schéma XML cible comprend dix champs par document : \texttt{<article>},
\texttt{<bulletin>}, \texttt{<date>}, \texttt{<rubrique>}, \texttt{<titre>},
\texttt{<auteur>}, \texttt{<texte>}, \texttt{<images>}, \texttt{<contact>}.
La figure~\ref{fig:architecture_td1} détaille les dépendances entre les modules
d'extraction.

\begin{figure}[ht]
  \centering
  \begin{tikzpicture}[
      mod/.style={draw, rounded corners=3pt, fill=codegray, minimum width=3cm,
                  minimum height=0.75cm, font=\small\ttfamily, align=center},
      arr/.style={-Stealth},
    ]
    \node[mod] (models)  {models.py};
    \node[mod, left=2cm of models]  (parser)  {parser.py};
    \node[mod, right=2cm of models] (builder) {xml\_builder.py};
    \node[mod, below=1.2cm of models] (pipeline) {pipeline.py};
    \draw[arr] (parser)  -- (models)   node[midway, above, font=\tiny] {Article};
    \draw[arr] (builder) -- (models)   node[midway, above, font=\tiny] {Article};
    \draw[arr] (pipeline) -- (parser)  node[near end, right, font=\tiny] {parse()};
    \draw[arr] (pipeline) -- (builder) node[near end, left,  font=\tiny] {add()};
  \end{tikzpicture}
  \caption{Dépendances entre les modules du TD1 : \texttt{models.py} est le seul point
    de contact entre parser et builder.}
  \label{fig:architecture_td1}
\end{figure}

La figure~\ref{fig:architecture_td1} montre que le module \texttt{models.py} joue le
rôle de contrat d'interface entre les deux modules principaux, sans qu'aucun n'ait
connaissance de l'implémentation de l'autre.

% ─────────────────────────────────────────────────────────────────────────────
\section{Architecture et conception}
% ─────────────────────────────────────────────────────────────────────────────

L'architecture suit une séparation claire des responsabilités. Le modèle de données
repose sur trois dataclasses : \texttt{Article} centralise tous les champs du corpus ;
\texttt{Person} structure l'auteur (nom, email) ; \texttt{Contact} enrichit les
informations \og{}Pour en savoir plus\fg{} avec quatre attributs distincts (nom, email,
URL, téléphone). Ce choix — dataclass plutôt que \texttt{dict} — permet d'avoir des
types annotés, une autocomplétion IDE, et de détecter dès la construction si un champ
est absent.

\begin{figure}[ht]
  \centering
  \begin{tikzpicture}[
      cls/.style={draw, rounded corners=2pt, fill=codegray,
                  inner sep=5pt, align=left, font=\scriptsize\ttfamily},
      comp/.style={-Stealth, thick},
      dep/.style={-Stealth, dashed, thick},
    ]
    % Nœud central
    \node[cls] (article) at (0,0) {
      \normalfont\textbf{Article}\\[2pt]
      code, bulletin : str\\
      date : date | None\\
      rubrique, title, body : str\\
      author : Person | None\\
      contacts : list[Contact]
    };
    % Dataclasses associées
    \node[cls] (person) at (6.5, 1.5) {
      \normalfont\textbf{Person}\\[2pt]
      name : str\\
      email : str
    };
    \node[cls] (contact) at (6.5, -1.5) {
      \normalfont\textbf{Contact}\\[2pt]
      name, email : str\\
      url, phone : str
    };
    % Classes de service
    \node[cls] (parser) at (-6, 1.5) {
      \normalfont\textbf{ArticleParser}\\[2pt]
      parse(path) : Article
    };
    \node[cls] (builder) at (-6, -1.5) {
      \normalfont\textbf{CorpusBuilder}\\[2pt]
      add\_article(a : Article)\\
      write(path : Path)
    };
    % Relations de composition
    \draw[comp] (article.east |- person.west) -- (person.west)
      node[midway, above, font=\tiny\normalfont] {author};
    \draw[comp] (article.east |- contact.west) -- (contact.west)
      node[midway, below, font=\tiny\normalfont] {contacts};
    % Relations de dépendance (crée / consomme)
    \draw[dep] (parser.east) -- (article.west |- parser.east)
      node[midway, above, font=\tiny\normalfont] {crée};
    \draw[dep] (builder.east) -- (article.west |- builder.east)
      node[midway, below, font=\tiny\normalfont] {consomme};
  \end{tikzpicture}
  \caption{Diagramme de classes du TD1 : \texttt{Article} est le pivot entre les classes
    de données (\texttt{Person}, \texttt{Contact}) et les classes de service
    (\texttt{ArticleParser}, \texttt{CorpusBuilder}).}
  \label{fig:classes_td1}
\end{figure}

La figure~\ref{fig:classes_td1} illustre que \texttt{ArticleParser} et
\texttt{CorpusBuilder} ne se connaissent pas mutuellement : \texttt{Article} est leur
seul point de couplage. Cette structure facilite les tests unitaires — chaque module
est testable indépendamment — et les évolutions futures, par exemple changer le format
de sortie sans toucher au parser.

Le flux d'exécution suit quatre phases séquentielles : (1) glob sur
\texttt{data/BULLETINS/*.htm} → 326 chemins ; (2) pour chaque fichier, parsing HTML
et construction d'un objet \texttt{Article} ; (3) vérification de doublon sur la clé
\texttt{(bulletin, code)} ; (4) sérialisation en \texttt{corpus.xml} puis rapport de
validation par champ.

% ─────────────────────────────────────────────────────────────────────────────
\section{Extraction champ par champ}
% ─────────────────────────────────────────────────────────────────────────────

Le module \texttt{parser.py} expose une fonction publique par champ. Chacune est testée
indépendamment avant d'être intégrée dans \texttt{ArticleParser.parse()}.

\paragraph{Numéro d'article (\texttt{parse\_code}).}
Recherche tous les \texttt{<span class="style15">} contenant le texte \og{}Code\fg{},
puis suit la balise \texttt{<a>} adjacente pour extraire l'identifiant numérique.

\begin{lstlisting}[caption={Extraction du numéro d'article par navigation sur les ancres},
  label={lst:parse_code}]
def parse_code(soup: BeautifulSoup) -> str:
    for span in soup.find_all("span", class_="style15"):
        if "Code" in span.get_text():
            a_tag = span.find_next("a")
            if isinstance(a_tag, Tag):
                return a_tag.get_text(strip=True)
    logger.warning("Missing article code")
    return ""
\end{lstlisting}

\paragraph{Numéro de bulletin (\texttt{parse\_bulletin}).}
Localise le \texttt{<span class="style32">} qui contient systématiquement l'identifiant
du bulletin (ex.\ \texttt{BE France 258}).

\paragraph{Rubrique (\texttt{parse\_rubrique}).}
Recherche le premier \texttt{<span class="style42">} à l'intérieur d'un
\texttt{<p class="style96">}. Ce sélecteur est distinct de celui utilisé pour la date,
car \texttt{style96} identifie les paragraphes du corps de l'article.

\paragraph{Titre (\texttt{parse\_title}).}
Localise l'unique \texttt{<span class="style17">} du fichier.

\paragraph{Auteur (\texttt{parse\_author}).}
Utilise le helper \texttt{\_find\_sibling\_cell()} pour localiser la cellule de tableau
voisine de l'étiquette \texttt{style28} « Rédacteur », puis extrait nom et email depuis
le \texttt{<span class="style95">} qu'elle contient. Le préfixe \og{}ADIT - \fg{} est
supprimé si présent.

\paragraph{Texte (\texttt{parse\_body}).}
Filtre les paragraphes \texttt{<p class="style96">} dont le premier enfant \texttt{<span>}
porte la classe \texttt{style95}. Ce filtrage distingue le corps textuel des rubriques et
dates qui partagent la même classe de paragraphe.

\paragraph{Images (\texttt{parse\_images}).}
Extrait les URLs des \texttt{<img>} présents dans la sidebar gauche
\texttt{<td class="FWExtra">}, en excluant les espaceurs (\texttt{\_clear.gif}) et les
icônes de navigation (\texttt{Resources/}). Les \texttt{<legendeImage>} sont
structurellement présentes dans le XML mais vides : les légendes ne sont pas encodées
dans le HTML source du corpus.

\paragraph{Contacts (\texttt{parse\_contacts}).}
Localise la cellule voisine de « Pour en savoir plus » puis parse chaque
\texttt{<p class="style44"> > <span class="style85">} par expressions régulières
pour extraire séparément nom, email (via \texttt{href="mailto:"} ou regex), URL et
téléphone. Les contacts sans nom, email ni URL sont filtrés.

\paragraph{Sérialisation XML (\texttt{CorpusBuilder.write}).}
La sérialisation finale utilise \texttt{lxml.etree} avec déclaration d'encodage UTF-8,
indentation et pretty-print. Le choix de \texttt{lxml.etree} (plutôt que la concaténation
de chaînes) garantit un XML valide et bien encodé.

\begin{lstlisting}[caption={Sérialisation finale du corpus en XML indenté UTF-8},
  label={lst:corpus_write}]
def write(self, output_path: Path) -> None:
    tree = etree.ElementTree(self._root)
    output_path.parent.mkdir(parents=True, exist_ok=True)
    tree.write(
        str(output_path),
        encoding="utf-8",
        xml_declaration=True,
        pretty_print=True,
    )
\end{lstlisting}

% ─────────────────────────────────────────────────────────────────────────────
\section{Extraction défensive des champs}
% ─────────────────────────────────────────────────────────────────────────────

La stratégie d'extraction est \textbf{défensive} : chaque extracteur retourne
\texttt{None} en cas d'absence du champ cible et journalise un avertissement, sans
jamais lever d'exception. Le corpus ADIT est bruité — des balises ou des classes CSS
peuvent être absentes dans une minorité de fichiers — et un crash sur un seul article
invaliderait l'ensemble du traitement.

\begin{lstlisting}[caption={Extraction de la date avec gestion défensive des cas manquants},
  label={lst:parse_date}]
def parse_date(soup: BeautifulSoup) -> date | None:
    style32 = soup.find("span", class_="style32")
    if style32 is None:
        logger.warning("Missing style32 span — date unavailable")
        return None
    # remonter au <p> parent pour trouver le style42 adjacent
    parent = style32.parent
    style42 = parent.find("span", class_="style42") if parent else None
    if style42 is None:
        return None
    raw = style42.get_text(strip=True)
    try:
        parts = raw.split("/")
        return date(int(parts[2]), int(parts[1]), int(parts[0]))
    except (ValueError, IndexError):
        logger.warning("Unparseable date value: %r", raw)
        return None
\end{lstlisting}

Le code~\ref{lst:parse_date} illustre le patron systématiquement appliqué. La navigation
dans les tableaux HTML pour l'auteur et les contacts repose sur un helper
\texttt{\_find\_sibling\_cell()} qui localise la cellule adjacente à une étiquette
de style \texttt{style28} donnée, évitant toute duplication de logique.

% ─────────────────────────────────────────────────────────────────────────────
\section{Résultats et couverture}
% ─────────────────────────────────────────────────────────────────────────────

Le pipeline produit 326 documents sans aucun doublon (vérification sur la clé
\texttt{(bulletin, code)}). Le tableau~\ref{tab:td1_coverage} présente la couverture
par champ à l'issue du traitement.

\begin{table}[ht]
\centering
\caption{Couverture du corpus après extraction (TD1, 326 articles)}
\label{tab:td1_coverage}
\begin{tabular}{lrrr}
\toprule
\textbf{Champ} & \textbf{Présent} & \textbf{Absent} & \textbf{Couverture} \\
\midrule
article  & 326 & 0 & 100\,\% \\
bulletin & 326 & 0 & 100\,\% \\
date     & 326 & 0 & 100\,\% \\
rubrique & 326 & 0 & 100\,\% \\
titre    & 326 & 0 & 100\,\% \\
auteur   & 326 & 0 & 100\,\% \\
texte    & 326 & 0 & 100\,\% \\
contact  & 325 & 1 &  99\,\% \\
images   & 326 & 0 & 100\,\% \\
\bottomrule
\end{tabular}
\end{table}

Le tableau~\ref{tab:td1_coverage} confirme une couverture quasi parfaite : seul un
article sur 326 est dépourvu de section \og{}Pour en savoir plus\fg{}. Les
\texttt{<legendeImage>} sont structurellement présentes dans le XML mais vides, les
légendes n'étant pas encodées dans le HTML source. Le corpus contient en outre
1\,630 entrées \texttt{<image>} au total, soit environ 5 images par article en moyenne.

La sortie du pipeline confirme l'absence de doublons et la couverture complète :

\begin{lstlisting}[language={}, basicstyle=\ttfamily\scriptsize,
  caption={Extrait du rapport de validation produit par \texttt{pipeline.py}},
  label={lst:pipeline_output}]
INFO: Found 326 bulletin files
INFO: Corpus written to outputs/corpus.xml (326 documents)
INFO: ====================================================
INFO: Corpus report — 326 articles written (0 duplicates skipped)
INFO: Field         Present  Missing   Coverage
INFO: code              326        0     100.0%
INFO: bulletin          326        0     100.0%
INFO: date              326        0     100.0%
INFO: rubrique          326        0     100.0%
INFO: title             326        0     100.0%
INFO: author            326        0     100.0%
INFO: body              326        0     100.0%
\end{lstlisting}

La suite de 75 tests unitaires s'exécute en 0,17\,s. Le tableau~\ref{tab:td1_coverage_code}
détaille la couverture par module.

\begin{table}[ht]
\centering
\caption{Couverture de code par module (TD1, pytest-cov)}
\label{tab:td1_coverage_code}
\begin{tabular}{lrrrl}
\toprule
\textbf{Module} & \textbf{Stmts} & \textbf{Miss} & \textbf{Cover} & \textbf{Lignes manquantes} \\
\midrule
\texttt{\_\_init\_\_.py}    &   0 &  0 & 100\,\% & — \\
\texttt{models.py}          &  23 &  0 & 100\,\% & — \\
\texttt{parser.py}          & 155 &  5 &  97\,\% & 58--59, 137, 217, 220 \\
\texttt{xml\_builder.py}    &  47 &  0 & 100\,\% & — \\
\midrule
\textbf{Total}              & 225 &  5 &  98\,\% & \\
\bottomrule
\end{tabular}
\end{table}

Le tableau~\ref{tab:td1_coverage_code} montre que les 5 lignes non couvertes dans
\texttt{parser.py} correspondent à des chemins d'erreur peu fréquents : l'avertissement
\og{}style32 has no parent \texttt{<p>}\fg{} (lignes 58--59), la construction d'URL sans
préfixe \texttt{http://} (ligne 137), et deux gardes défensives dans \texttt{parse\_contacts}
pour des éléments non-\texttt{Tag} (lignes 217, 220).

% ─────────────────────────────────────────────────────────────────────────────
\section{Difficultés rencontrées}
% ─────────────────────────────────────────────────────────────────────────────

\paragraph{Encodage des fichiers HTML.}
Les fichiers du corpus déclarent \texttt{charset=ISO-8859-1} dans leur en-tête HTML,
mais le sujet précise qu'ils sont en UTF-8. L'utilisation de \texttt{path.read\_bytes()}
plutôt que \texttt{open(\ldots, encoding=\ldots)} permet de laisser BeautifulSoup détecter
l'encodage correct automatiquement via les métadonnées du document.

\paragraph{Distinction rubrique / corps de l'article.}
Les paragraphes de rubrique et les paragraphes de texte partagent la classe \texttt{style96}.
La distinction repose sur la classe du premier enfant \texttt{<span>} : \texttt{style42}
pour les métadonnées (rubrique, date) et \texttt{style95} pour le texte. Ce comportement
n'est pas documenté dans les fichiers HTML et a nécessité une exploration manuelle de
plusieurs bulletins.

\paragraph{Navigation par tableaux pour l'auteur et les contacts.}
Les champs « Rédacteur » et « Pour en savoir plus » ne sont pas identifiés par des
attributs stables, mais par le texte d'une cellule de tableau voisine. Le helper
\texttt{\_find\_sibling\_cell()} centralise cette logique de navigation pour les deux
champs, évitant toute duplication de code.

% ─────────────────────────────────────────────────────────────────────────────
\section{Fonctionnalités supplémentaires}
% ─────────────────────────────────────────────────────────────────────────────

\paragraph{Modélisation riche des contacts.}
Au-delà de la consigne minimale (un champ \texttt{<contact>} textuel), les contacts
sont modélisés par un dataclass \texttt{Contact} avec quatre attributs distincts
(\texttt{name}, \texttt{email}, \texttt{url}, \texttt{phone}). Ce choix anticipe les
besoins des TD suivants : l'indexation séparée des adresses email et des URLs de contact
est beaucoup plus facile à partir de données structurées que depuis une chaîne brute.

\paragraph{Mode échantillon (\texttt{SAMPLE\_SIZE}).}
Le pipeline accepte la variable d'environnement \texttt{SAMPLE\_SIZE} pour limiter le
traitement à $N$ fichiers, accélérant le cycle de développement lors des tests manuels :
\texttt{SAMPLE\_SIZE=10 uv run python -m adit\_corpus\_indexing.pipeline}.

\paragraph{Déduplication des articles.}
Le pipeline détecte et ignore les doublons sur la clé \texttt{(bulletin, code)} en
journalisant un avertissement. Aucun doublon n'a été trouvé dans le corpus actuel, mais
ce mécanisme assure la robustesse si le répertoire venait à contenir des copies.
